% Build from the repository root:
% latexmk -pdf -interaction=nonstopmode -halt-on-error -outdir=tmp/section34_spectrum docs/section34_spectrum_note.tex
% The proposed paper text is in section34_spectrum_subsection.tex.
\documentclass[10pt,letterpaper]{article}
\usepackage[textwidth=5.5in,textheight=9in,centering]{geometry}
\usepackage{times,amsmath,amssymb,amsthm,graphicx,microtype,flafter,placeins,xurl}
\usepackage[round,authoryear]{natbib}
\usepackage[colorlinks=true,linkcolor=blue,urlcolor=blue,citecolor=blue]{hyperref}
\usepackage{titlesec}
\titleformat{\subsection}{\normalfont\normalsize\scshape}{\thesubsection}{1em}{}
\setlength{\parindent}{0pt}
\setlength{\parskip}{5pt plus 1pt minus 1pt}
\widowpenalty=10000
\clubpenalty=10000
\newtheorem{theorem}{Theorem}[section]
\newtheorem{corollary}[theorem]{Corollary}
\DeclareMathOperator{\sech}{sech}
\DeclareMathOperator{\sign}{sign}
\title{Bandwidth controls optimization access to high precision}
\author{Section 3.4 working note}
\date{September 24, 2026}
\begin{document}
\maketitle
\vspace{-1.3em}
\noindent\textit{Proposed subsection, followed by its full spectral argument and
evaluation details. References to Sections 3.1--3.3 and Figure 2b refer to
submission draft (19). The final paragraph leads into the forthcoming
joint-feature analysis.}

\setcounter{section}{3}
\setcounter{subsection}{3}
\setcounter{theorem}{1}
\setcounter{figure}{2}
\input{docs/section34_spectrum_subsection.tex}
\FloatBarrier
\clearpage
\appendix
\input{docs/section34_spectrum_appendix.tex}
\bibliographystyle{plainnat}
\bibliography{docs/section34_references}
\end{document}
